\section{Experimental design}\label{sec:design}

The experiments ask whether the front computed by continuation is worth its cost against the points that the published methods give on the same array, and they measure the properties that the theory predicts. Every method is run by the engine of the studio on the same task and seeds, and the values reported in the sources are used only as a check of the order of
magnitude of the published points.

\subsection{Hypotheses}\label{sec:hypotheses}

The hypotheses follow from the acceptance criterion of the method and from the propositions of Section~\ref{sec:theory}; they are stated here before the results and each has a rule that decides it.
Table~\ref{tab:hyp} lists them. The first three come from the acceptance criterion of the module, which asks that the front have a hypervolume at least that of the union of the points of Gerchberg--Saxton, the adaptive additive iteration, weighted GS and direct search, that at least one front point dominate each of the first two, and that the front take less than 60 seconds; the research record
that motivated the method named a risk, that the point of weighted GS is already near the best attainable corner, so that the front would be flat, and the fourth and fifth hypotheses test it. The remaining ones test the theory and the design choices of the method.

\begin{table}[t]
\centering\small
\caption{Hypotheses, the measure that tests each, and the rule that decides it. Seeds are the seeds of the main experiment unless a hypothesis names the ablation seeds; all intervals are 95 percent intervals of the mean over seeds.}
\label{tab:hyp}
\begin{tabular}{@{}p{0.7cm} p{6.0cm} p{7.4cm}@{}}
\toprule
 & Hypothesis & Decision rule\\
\midrule
H1 & The hypervolume of the front is at least that of the union of the points of Gerchberg--Saxton, adaptive additive, weighted GS and direct search & the paired difference per seed has an interval above zero and no seed is negative\\
H2 & In every seed at least one front point dominates the point of Gerchberg--Saxton and one the point of the adaptive additive iteration & the count of dominating points is at least one for both points in every seed\\
H3 & The front completes in under 60 seconds on the CPU & the time of the module's own run, alone, under 60 s in every timed seed\\
H4 & At least one front point dominates the point of weighted GS, of phase-fixed WGS and of direct search in every seed & the count of dominating points is at least one in every seed for each of the three\\
H5 & The front has a hypervolume at least that of the union of the one-parameter fronts that the published methods give when their iteration count or exponent is varied & the paired difference per seed has an interval above zero\\
H6 & The continuation is monotone and obeys the bounds: the efficiency and the error $E$ do not decrease along $\lambda$, $\sigma$ and $u$ respect Proposition~\ref{prop:bound} & no violation in any seed and no negative slack on any front point\\
H7 & The warm start is worth its design: the cold start at equal cost gives a lower hypervolume & the paired difference of the ablation has an interval below zero\\
H8 & The stored fronts of the benchmark meet the hypervolume of the published union on every geometry & the margin of each geometry, at 30 and at 100 iterations, is not negative\\
\bottomrule
\end{tabular}
\end{table}

\subsection{Data and provenance}\label{sec:data}

Three sources of data are used, and each is read or generated by a committed script. The first is the acceptance task defined in Section~\ref{sec:problem}, on which the front and the published points are computed afresh for seeds $1,\dots,\nSeeds$ (the main experiment) and on which the ablation runs for the seeds $1,\dots,\nSeedsAbl$. The second is the benchmark of the studio, a table of \benchRows{} rows of the engine at release \benchEngine{} over seven geometries,
three iteration budgets and two wall-clock budgets, with medians over up to \benchSeeds{} seeds of each cell, and a stored front for each geometry that the method was run on. The geometries are a ten by ten lattice, a ring of 36 traps, 100 random traps, a quasi-periodic set of 119 traps, a three-dimensional lattice of $5\times5\times3$ traps and a set of 1024 traps; they are read from the file and not re-run.
The third is the bake of the studio's cases, in which each case is run with the method for three seeds and the knee is stored with its efficiency, uniformity and spread; it is read from the committed file of measured ranges. The first source tests the hypotheses on the task of the foundational comparison, the second tests whether the result carries to other geometries, and the third gives the
spread of the knee that the studio shows.

\subsection{Methods compared}\label{sec:compared}

The points of the published methods are computed by the engine on the same task and seeds as the front, never copied from the sources. The engine runs the rule of each method through its own loop; the start of every method is the same seeded random superposition. Table~\ref{tab:methods} lists the methods, their settings and their origin.

\begin{table}[t]
\centering\small
\caption{Published methods that the engine runs on the task as points of the plane $(e,u)$, with their settings. The iteration count of the first four is that of the acceptance test of the method; the settings of direct search are the engine defaults.}
\label{tab:methods}
\begin{tabular}{@{}p{4.2cm} p{4.4cm} p{5.4cm}@{}}
\toprule
Method & Setting & Origin of the rule\\
\midrule
Gerchberg--Saxton & \baseIters{} iterations & Di Leonardo et al.~\cite{dileonardo}\\
Adaptive additive & \baseIters{} iterations, exponent 0.5 & Di Leonardo et al.~\cite{dileonardo}\\
Weighted GS & \baseIters{} iterations & Di Leonardo et al.~\cite{dileonardo}\\
Phase-fixed WGS & \baseIters{} iterations, freezing at 12 & Kim et al.~\cite{kim2019}\\
Direct search & one sweep, 8 trial phases, weight $f=0.5$ & Polin et al.~\cite{polin}, Di Leonardo et al.~\cite{dileonardo}\\
\bottomrule
\end{tabular}
\end{table}

The first four rows are the points of the acceptance criterion of the method together with weighted GS; phase-fixed WGS is added here because it is the start of the front, so that its point is the uniform end of every front, up to the refinement by the first minimization. The union of the points of Gerchberg--Saxton, the adaptive additive iteration, weighted GS and direct search is the published union of hypotheses H1 and H8.
For hypothesis H5 the engine also generates three one-parameter fronts from the published rules, each a set of points that a reader of the sources could produce by turning the one knob that the source exposes: the iterates $1,\dots,100$ of weighted GS, the iterates $1,\dots,100$ of phase-fixed WGS, and the adaptive additive iteration at \baseIters{} iterations with its exponent set to $0,0.1,\dots,1$. The sweep of the weight of direct search is not run, because one sweep of one value takes about a minute on the CPU of the machine.

\subsection{Metrics and statistics}\label{sec:metrics}

The metrics are the efficiency and the uniformity of Eq.~\ref{eq:eu}, the relative standard deviation of Eq.~\ref{eq:sigma}, the hypervolume of Eq.~\ref{eq:hv} with reference $(0,0)$ and with two further references for the sensitivity analysis, the number of front points that dominate a published point, and the wall-clock seconds. The dominance count of a point $(e_b,u_b)$ is
\begin{equation}
N_{\mathrm{dom}}=\big|\{k:\ e_k\ge e_b,\ u_k>u_b\}\big|,\label{eq:ndom}
\end{equation}
the criterion of the module's own test. For $n$ seeds a statistic $x_i$ is summarized by its mean, its sample standard deviation and the two-sided 95 percent interval of the mean,
\begin{equation}
\bar x=\frac1n\sum_ix_i,\qquad s^2=\frac1{n-1}\sum_i(x_i-\bar x)^2,\qquad \bar x\pm t_{n-1}\frac{s}{\sqrt n},\label{eq:ci}
\end{equation}
with $t_{n-1}$ the quantile of Student's distribution, and a comparison of two methods is the paired difference over the seeds, $\Delta_i=\mathrm{HV}^{\mathrm{front}}_i-\mathrm{HV}^{\mathrm{ref}}_i$, with the interval of Eq.~\ref{eq:ci} applied to $\Delta_i$. The pairing is by seed: both methods start from the same random superposition, so the difference removes the variation that comes from the start.
With eight seeds the quantile is 2.365 and with five 2.776, so the intervals are wide and the tests that they imply have low power; no $p$-value is reported and a hypothesis is decided by the rule of Table~\ref{tab:hyp}, with the interval and the number of seeds on each side of zero stated. The counts of seeds that satisfy a rule are reported next to the interval so that a reader can see the direction of every seed.

\subsection{Protocol and ablations}\label{sec:protocol}

Every method is deterministic given its seed through the portable generator of the engine, so a rerun reproduces the numbers up to the floating point of the machine's transform. The main experiment runs the front and the five points for $\nSeeds$ seeds, and the one-parameter fronts for the same seeds. The time of the front is measured separately, alone on the machine, with the module's own entry and the threads of the machine, for three seeds,
because the runs of the main experiment execute in parallel processes with one thread each and their seconds are not the time of the method. The loop that the ablation uses is a copy of the module's loop with the choices that it varies; the script checks that for the default choices it reproduces the front of the module for the first seed exactly.

The ablation varies one choice at a time from the reference, on the first $\nSeedsAbl$ seeds. The number of weights is 5, 10 and 40 against 20. The conjugate-gradient iterations per weight are 5 and 30 against 15. The schedule of the weights is quadratic in the grid index, which is dense near $\lambda=0$, and target-spaced, which is linear in the requested total light and ends at the largest
total that the optics allow (Section~\ref{sec:theory}). The continuation is replaced by a cold start, in which every weight starts from the refined hologram of the start and not from its neighbour, at the same number of steps per weight. The start is replaced by 30 iterations of plain GS, by the random superposition with no iteration, and by weighted GS after 100 iterations. Finally the loss is replaced by the unscaled form of the research record,
$(1-\lambda)\sum_m(I_m-T_m)^2-\lambda\sum_mI_m$, which tests the normalization. Each variant reports the hypervolume, its paired difference from the reference, the lowest uniformity reached, the largest efficiency and the seconds of one thread.

A diagnostic separates the explanations of a result that the main experiment produces, the domination of part of the front by iterates of weighted GS. For each seed it scores the holograms of weighted GS after 30 and 100 iterations and of phase-fixed WGS after 100 iterations on the normalized error $E$ and the normalized light $S$ of that seed's target, computes their loss $L_\lambda$ at each weight of the seed's front, and counts the weights at which the published hologram has a lower loss
than the front point. A positive count means that the minimization of the front stopped at a point that is worse, on the loss it minimizes, than a point that the published method reaches. The diagnostic uses the same target total $e_0$ as the front, so the comparison is on one loss.
